% ── Pokrytí KV – časový vývoj (cb_coverage_timeline.py) ─────────────────────

Dlouhodobý pohled na vývoj pokrytí \ac{KS} sahající do počátku 90.~let
(obr.~\ref{fig:cb_coverage_timeline}) ukazuje, že bimodální rozdělení
evropských zemí na skupinu s~vysokým pokrytím (nordické a~kontinentální
státy) a~skupinu s~nízkým pokrytím (středo- a~východoevropské státy) není
důsledkem transformačního přechodu devadesátých let ani vstupu do \ac{EU},
nýbrž stav hluboce zakořeněný v~institucionální architektuře \ac{KV}.
Rakousko~(AT) udržuje pokrytí blízké~\SI{98}{\percent} po celé sledované
období díky zákonné autoritě sektorových hospodářských komor a~automatickému
rozšiřování závaznosti \ac{KS}; Dánsko~(DK) a~ostatní severské země se
pohybují stabilně mezi \SI{80}{\percent} a~\SI{90}{\percent} navzdory klesající
hustotě odborů --- mechanismem erga omnes zajišťují pokrytí i~pracovníkům
mimo odbory.
Naproti tomu ČR a~Polsko~(PL) vykazovaly nízké a~klesající pokrytí
\emph{ještě před vstupem do \ac{EU}} --- institucionální deficit tedy
nepředstavuje výsledek liberalizace po~2004, ale trvalé charakteristiky
systému, který nikdy nevyvinul sektorové erga omnes.

% ── Pokrytí KV – časový vývoj (cb_coverage_timeline.py) ─────────────────────

Dlouhodobý pohled na vývoj pokrytí \ac{KS} sahající do počátku 90.~let
(obr.~\ref{fig:cb_coverage_timeline}) ukazuje, že bimodální rozdělení
evropských zemí na skupinu s~vysokým pokrytím (nordické a~kontinentální
státy) a~skupinu s~nízkým pokrytím (středo- a~východoevropské státy) není
důsledkem transformačního přechodu devadesátých let ani vstupu do \ac{EU},
nýbrž stav hluboce zakořeněný v~institucionální architektuře \ac{KV}.
Rakousko~(AT) udržuje pokrytí blízké~\SI{98}{\percent} po celé sledované
období díky zákonné autoritě sektorových hospodářských komor a~automatickému
rozšiřování závaznosti \ac{KS}; Dánsko~(DK) a~ostatní severské země se
pohybují stabilně mezi \SI{80}{\percent} a~\SI{90}{\percent} navzdory klesající
hustotě odborů --- mechanismem erga omnes zajišťují pokrytí i~pracovníkům
mimo odbory.
Naproti tomu ČR a~Polsko~(PL) vykazovaly nízké a~klesající pokrytí
\emph{ještě před vstupem do \ac{EU}} --- institucionální deficit tedy
nepředstavuje výsledek liberalizace po~2004, ale trvalé charakteristiky
systému, který nikdy nevyvinul sektorové erga omnes.

% ── Pokrytí KV – časový vývoj (cb_coverage_timeline.py) ─────────────────────

Dlouhodobý pohled na vývoj pokrytí \ac{KS} sahající do počátku 90.~let
(obr.~\ref{fig:cb_coverage_timeline}) ukazuje, že bimodální rozdělení
evropských zemí na skupinu s~vysokým pokrytím (nordické a~kontinentální
státy) a~skupinu s~nízkým pokrytím (středo- a~východoevropské státy) není
důsledkem transformačního přechodu devadesátých let ani vstupu do \ac{EU},
nýbrž stav hluboce zakořeněný v~institucionální architektuře \ac{KV}.
Rakousko~(AT) udržuje pokrytí blízké~\SI{98}{\percent} po celé sledované
období díky zákonné autoritě sektorových hospodářských komor a~automatickému
rozšiřování závaznosti \ac{KS}; Dánsko~(DK) a~ostatní severské země se
pohybují stabilně mezi \SI{80}{\percent} a~\SI{90}{\percent} navzdory klesající
hustotě odborů --- mechanismem erga omnes zajišťují pokrytí i~pracovníkům
mimo odbory.
Naproti tomu ČR a~Polsko~(PL) vykazovaly nízké a~klesající pokrytí
\emph{ještě před vstupem do \ac{EU}} --- institucionální deficit tedy
nepředstavuje výsledek liberalizace po~2004, ale trvalé charakteristiky
systému, který nikdy nevyvinul sektorové erga omnes.

% ── Pokrytí KV – časový vývoj (cb_coverage_timeline.py) ─────────────────────

Dlouhodobý pohled na vývoj pokrytí \ac{KS} sahající do počátku 90.~let
(obr.~\ref{fig:cb_coverage_timeline}) ukazuje, že bimodální rozdělení
evropských zemí na skupinu s~vysokým pokrytím (nordické a~kontinentální
státy) a~skupinu s~nízkým pokrytím (středo- a~východoevropské státy) není
důsledkem transformačního přechodu devadesátých let ani vstupu do \ac{EU},
nýbrž stav hluboce zakořeněný v~institucionální architektuře \ac{KV}.
Rakousko~(AT) udržuje pokrytí blízké~\SI{98}{\percent} po celé sledované
období díky zákonné autoritě sektorových hospodářských komor a~automatickému
rozšiřování závaznosti \ac{KS}; Dánsko~(DK) a~ostatní severské země se
pohybují stabilně mezi \SI{80}{\percent} a~\SI{90}{\percent} navzdory klesající
hustotě odborů --- mechanismem erga omnes zajišťují pokrytí i~pracovníkům
mimo odbory.
Naproti tomu ČR a~Polsko~(PL) vykazovaly nízké a~klesající pokrytí
\emph{ještě před vstupem do \ac{EU}} --- institucionální deficit tedy
nepředstavuje výsledek liberalizace po~2004, ale trvalé charakteristiky
systému, který nikdy nevyvinul sektorové erga omnes.

\input{texparts/python/cb_coverage_timeline}

Přiblížení na období 2004--2024 (obr.~\ref{fig:cb_coverage_timeline_2004})
odhaluje diferenciaci trajektorií v~rámci šedého oblaku \ac{EU}27.
Šedé linie nejsou homogenní: přibližně třetina zemí si zachovává pokrytí
nad \SI{70}{\percent} (FR, BE, IT, NL a~nordické státy), zbývající dvě třetiny
se pohybují pod \SI{50}{\percent} a~tvoří sestupnou skupinu.
ČR zaujímá v~tomto kontextu výraznou pozici --- s~pokrytím přibližně
\SI{31}{\percent} (2023) se nachází těsně nad skupinou nejméně pokrytých zemí
a~její trajektorie stagnuje bez signálu obnovy.
Slovensko~(SK) vykazuje podobnou hodnotu, nicméně s~méně výrazným poklesem
v~prvním desetiletí po vstupu do \ac{EU}.
Německo~(DE) s~hodnotou ERB přibližně \SI{49}{\percent} (2024) ilustruje
zranitelnost i~u~ekonomik s~dlouhou tradicí \ac{KV}: postupný ústup
odvětvového vyjednávání ve prospěch podnikové úrovně od 90.~let snižoval
pokrytí; hodnota ERB přitom nadhodnocuje skutečný dosah kolektivně dohodnutých
podmínek, protože ERB zachycuje registrované smlouvy, nikoliv podíl
zaměstnanců fakticky pokrytých (metodický rozdíl vůči \textit{AdjCov} ostatních
zemí je uveden v~záhlaví obrázku).
Polsko~(PL) se celé období drží hluboko pod \SI{20}{\percent} a~v~posledních
letech stagnuje nejníže ze zobrazených zemí.

\input{texparts/python/cb_coverage_timeline_2004}

Porovnání obr.~\ref{fig:cb_coverage_timeline_2004} s~vývojem hustoty odborů
(obr.~\ref{fig:union_density_trend_2004}) odhaluje klíčový strukturní poznatek:
\emph{pokrytí a~hustota se vyvíjejí nezávisle}.
Francie kombinuje hustotu odborů pod \SI{8}{\percent} s~pokrytím přesahujícím
\SI{90}{\percent} díky generálnímu rozšiřování závaznosti sektorových smluv.
Naopak ČR vykazuje souběžný pokles obou ukazatelů, neboť bez erga omnes
mechanismu se pokrytí odvíjí přímo od počtu uzavřených podnikových smluv,
jejichž dosah se s~fragmentací pracovního trhu zužuje.
Skutečnost, že severské země udržují vysoké pokrytí i~při relativně klesající
hustotě odborů, svědčí o~tom, že institucionální rámec erga omnes je
stabilizačním faktorem oddělujícím pokrytí od členství.
Reforma \ac{KV} v~ČR zaměřená \emph{pouze} na nábor odborových členů by proto
byla méně účinná než systémová změna mechanismu závaznosti.

Trajektorie pokrytí \ac{KS} zobrazená v~obr.~\ref{fig:cb_coverage_timeline_2004}
zároveň tvoří ústřední nezávislou proměnnou korelační analýzy
(odst.~\ref{ssec:korelace}): pozice ČR na přibližně \SI{31}{\percent} ji
konzistentně umísťuje do kvadrantu horší výkonnosti ve všech čtyřech
sledovaných párování --- s~delší pracovní dobou, nižším HDP/ob.~v~PPS,
vyšším Giniho koeficientem a~nižším čistým příjmem v~PPS.
Modelový výpočet v~odst.~\ref{par:model_cz} kvantifikuje, jakou změnu
výstupních ukazatelů by přineslo přiblížení k~cílové hodnotě~\SI{80}{\percent}
dle směrnice~2022/2041.


Přiblížení na období 2004--2024 (obr.~\ref{fig:cb_coverage_timeline_2004})
odhaluje diferenciaci trajektorií v~rámci šedého oblaku \ac{EU}27.
Šedé linie nejsou homogenní: přibližně třetina zemí si zachovává pokrytí
nad \SI{70}{\percent} (FR, BE, IT, NL a~nordické státy), zbývající dvě třetiny
se pohybují pod \SI{50}{\percent} a~tvoří sestupnou skupinu.
ČR zaujímá v~tomto kontextu výraznou pozici --- s~pokrytím přibližně
\SI{31}{\percent} (2023) se nachází těsně nad skupinou nejméně pokrytých zemí
a~její trajektorie stagnuje bez signálu obnovy.
Slovensko~(SK) vykazuje podobnou hodnotu, nicméně s~méně výrazným poklesem
v~prvním desetiletí po vstupu do \ac{EU}.
Německo~(DE) s~hodnotou ERB přibližně \SI{49}{\percent} (2024) ilustruje
zranitelnost i~u~ekonomik s~dlouhou tradicí \ac{KV}: postupný ústup
odvětvového vyjednávání ve prospěch podnikové úrovně od 90.~let snižoval
pokrytí; hodnota ERB přitom nadhodnocuje skutečný dosah kolektivně dohodnutých
podmínek, protože ERB zachycuje registrované smlouvy, nikoliv podíl
zaměstnanců fakticky pokrytých (metodický rozdíl vůči \textit{AdjCov} ostatních
zemí je uveden v~záhlaví obrázku).
Polsko~(PL) se celé období drží hluboko pod \SI{20}{\percent} a~v~posledních
letech stagnuje nejníže ze zobrazených zemí.

\input{texparts/python/cb_coverage_timeline_2004}

Porovnání obr.~\ref{fig:cb_coverage_timeline_2004} s~vývojem hustoty odborů
(obr.~\ref{fig:union_density_trend_2004}) odhaluje klíčový strukturní poznatek:
\emph{pokrytí a~hustota se vyvíjejí nezávisle}.
Francie kombinuje hustotu odborů pod \SI{8}{\percent} s~pokrytím přesahujícím
\SI{90}{\percent} díky generálnímu rozšiřování závaznosti sektorových smluv.
Naopak ČR vykazuje souběžný pokles obou ukazatelů, neboť bez erga omnes
mechanismu se pokrytí odvíjí přímo od počtu uzavřených podnikových smluv,
jejichž dosah se s~fragmentací pracovního trhu zužuje.
Skutečnost, že severské země udržují vysoké pokrytí i~při relativně klesající
hustotě odborů, svědčí o~tom, že institucionální rámec erga omnes je
stabilizačním faktorem oddělujícím pokrytí od členství.
Reforma \ac{KV} v~ČR zaměřená \emph{pouze} na nábor odborových členů by proto
byla méně účinná než systémová změna mechanismu závaznosti.

Trajektorie pokrytí \ac{KS} zobrazená v~obr.~\ref{fig:cb_coverage_timeline_2004}
zároveň tvoří ústřední nezávislou proměnnou korelační analýzy
(odst.~\ref{ssec:korelace}): pozice ČR na přibližně \SI{31}{\percent} ji
konzistentně umísťuje do kvadrantu horší výkonnosti ve všech čtyřech
sledovaných párování --- s~delší pracovní dobou, nižším HDP/ob.~v~PPS,
vyšším Giniho koeficientem a~nižším čistým příjmem v~PPS.
Modelový výpočet v~odst.~\ref{par:model_cz} kvantifikuje, jakou změnu
výstupních ukazatelů by přineslo přiblížení k~cílové hodnotě~\SI{80}{\percent}
dle směrnice~2022/2041.


Přiblížení na období 2004--2024 (obr.~\ref{fig:cb_coverage_timeline_2004})
odhaluje diferenciaci trajektorií v~rámci šedého oblaku \ac{EU}27.
Šedé linie nejsou homogenní: přibližně třetina zemí si zachovává pokrytí
nad \SI{70}{\percent} (FR, BE, IT, NL a~nordické státy), zbývající dvě třetiny
se pohybují pod \SI{50}{\percent} a~tvoří sestupnou skupinu.
ČR zaujímá v~tomto kontextu výraznou pozici --- s~pokrytím přibližně
\SI{31}{\percent} (2023) se nachází těsně nad skupinou nejméně pokrytých zemí
a~její trajektorie stagnuje bez signálu obnovy.
Slovensko~(SK) vykazuje podobnou hodnotu, nicméně s~méně výrazným poklesem
v~prvním desetiletí po vstupu do \ac{EU}.
Německo~(DE) s~hodnotou ERB přibližně \SI{49}{\percent} (2024) ilustruje
zranitelnost i~u~ekonomik s~dlouhou tradicí \ac{KV}: postupný ústup
odvětvového vyjednávání ve prospěch podnikové úrovně od 90.~let snižoval
pokrytí; hodnota ERB přitom nadhodnocuje skutečný dosah kolektivně dohodnutých
podmínek, protože ERB zachycuje registrované smlouvy, nikoliv podíl
zaměstnanců fakticky pokrytých (metodický rozdíl vůči \textit{AdjCov} ostatních
zemí je uveden v~záhlaví obrázku).
Polsko~(PL) se celé období drží hluboko pod \SI{20}{\percent} a~v~posledních
letech stagnuje nejníže ze zobrazených zemí.

\input{texparts/python/cb_coverage_timeline_2004}

Porovnání obr.~\ref{fig:cb_coverage_timeline_2004} s~vývojem hustoty odborů
(obr.~\ref{fig:union_density_trend_2004}) odhaluje klíčový strukturní poznatek:
\emph{pokrytí a~hustota se vyvíjejí nezávisle}.
Francie kombinuje hustotu odborů pod \SI{8}{\percent} s~pokrytím přesahujícím
\SI{90}{\percent} díky generálnímu rozšiřování závaznosti sektorových smluv.
Naopak ČR vykazuje souběžný pokles obou ukazatelů, neboť bez erga omnes
mechanismu se pokrytí odvíjí přímo od počtu uzavřených podnikových smluv,
jejichž dosah se s~fragmentací pracovního trhu zužuje.
Skutečnost, že severské země udržují vysoké pokrytí i~při relativně klesající
hustotě odborů, svědčí o~tom, že institucionální rámec erga omnes je
stabilizačním faktorem oddělujícím pokrytí od členství.
Reforma \ac{KV} v~ČR zaměřená \emph{pouze} na nábor odborových členů by proto
byla méně účinná než systémová změna mechanismu závaznosti.

Trajektorie pokrytí \ac{KS} zobrazená v~obr.~\ref{fig:cb_coverage_timeline_2004}
zároveň tvoří ústřední nezávislou proměnnou korelační analýzy
(odst.~\ref{ssec:korelace}): pozice ČR na přibližně \SI{31}{\percent} ji
konzistentně umísťuje do kvadrantu horší výkonnosti ve všech čtyřech
sledovaných párování --- s~delší pracovní dobou, nižším HDP/ob.~v~PPS,
vyšším Giniho koeficientem a~nižším čistým příjmem v~PPS.
Modelový výpočet v~odst.~\ref{par:model_cz} kvantifikuje, jakou změnu
výstupních ukazatelů by přineslo přiblížení k~cílové hodnotě~\SI{80}{\percent}
dle směrnice~2022/2041.


Přiblížení na období 2004--2024 (obr.~\ref{fig:cb_coverage_timeline_2004})
odhaluje diferenciaci trajektorií v~rámci šedého oblaku \ac{EU}27.
Šedé linie nejsou homogenní: přibližně třetina zemí si zachovává pokrytí
nad \SI{70}{\percent} (FR, BE, IT, NL a~nordické státy), zbývající dvě třetiny
se pohybují pod \SI{50}{\percent} a~tvoří sestupnou skupinu.
ČR zaujímá v~tomto kontextu výraznou pozici --- s~pokrytím přibližně
\SI{31}{\percent} (2023) se nachází těsně nad skupinou nejméně pokrytých zemí
a~její trajektorie stagnuje bez signálu obnovy.
Slovensko~(SK) vykazuje podobnou hodnotu, nicméně s~méně výrazným poklesem
v~prvním desetiletí po vstupu do \ac{EU}.
Německo~(DE) s~hodnotou ERB přibližně \SI{49}{\percent} (2024) ilustruje
zranitelnost i~u~ekonomik s~dlouhou tradicí \ac{KV}: postupný ústup
odvětvového vyjednávání ve prospěch podnikové úrovně od 90.~let snižoval
pokrytí; hodnota ERB přitom nadhodnocuje skutečný dosah kolektivně dohodnutých
podmínek, protože ERB zachycuje registrované smlouvy, nikoliv podíl
zaměstnanců fakticky pokrytých (metodický rozdíl vůči \textit{AdjCov} ostatních
zemí je uveden v~záhlaví obrázku).
Polsko~(PL) se celé období drží hluboko pod \SI{20}{\percent} a~v~posledních
letech stagnuje nejníže ze zobrazených zemí.

\input{texparts/python/cb_coverage_timeline_2004}

Porovnání obr.~\ref{fig:cb_coverage_timeline_2004} s~vývojem hustoty odborů
(obr.~\ref{fig:union_density_trend_2004}) odhaluje klíčový strukturní poznatek:
\emph{pokrytí a~hustota se vyvíjejí nezávisle}.
Francie kombinuje hustotu odborů pod \SI{8}{\percent} s~pokrytím přesahujícím
\SI{90}{\percent} díky generálnímu rozšiřování závaznosti sektorových smluv.
Naopak ČR vykazuje souběžný pokles obou ukazatelů, neboť bez erga omnes
mechanismu se pokrytí odvíjí přímo od počtu uzavřených podnikových smluv,
jejichž dosah se s~fragmentací pracovního trhu zužuje.
Skutečnost, že severské země udržují vysoké pokrytí i~při relativně klesající
hustotě odborů, svědčí o~tom, že institucionální rámec erga omnes je
stabilizačním faktorem oddělujícím pokrytí od členství.
Reforma \ac{KV} v~ČR zaměřená \emph{pouze} na nábor odborových členů by proto
byla méně účinná než systémová změna mechanismu závaznosti.

Trajektorie pokrytí \ac{KS} zobrazená v~obr.~\ref{fig:cb_coverage_timeline_2004}
zároveň tvoří ústřední nezávislou proměnnou korelační analýzy
(odst.~\ref{ssec:korelace}): pozice ČR na přibližně \SI{31}{\percent} ji
konzistentně umísťuje do kvadrantu horší výkonnosti ve všech čtyřech
sledovaných párování --- s~delší pracovní dobou, nižším HDP/ob.~v~PPS,
vyšším Giniho koeficientem a~nižším čistým příjmem v~PPS.
Modelový výpočet v~odst.~\ref{par:model_cz} kvantifikuje, jakou změnu
výstupních ukazatelů by přineslo přiblížení k~cílové hodnotě~\SI{80}{\percent}
dle směrnice~2022/2041.
